
\section{Compatibility target: factoring the Langlands upgrade through periods and motives}
\label{sec:langlands_compatibility}

The earlier ``stairway'' work \cite{Ma2025STI} emphasized a functorial upgrade objective: to construct a holographic Langlands functor from a category of scan protocols into a category of automorphic representations, with provable compatibility across renormalization flow, cusp discretization, and Hecke dynamics. The present paper contributes a new, auditable arrow---the functor $\HSP:\ScanAlg\to\PerDatum$---and uses it to formulate a refined compatibility target.

\subsection{A commutative-diagram objective}
\label{sec:commutative_diagram_objective}

In the classical modular setting, three layers are tightly related:
\begin{itemize}
\item \textbf{Period layer}: integrals of algebraic or modular differential forms against cycles, yielding periods (Section~\ref{sec:kz_periods}).
\item \textbf{$\ell$-adic layer}: Galois representations whose Frobenius traces encode arithmetic spectra, such as Hecke eigenvalues in the modular case \cite{Deligne1971Bourbaki,Deligne1974}.
\item \textbf{Automorphic layer}: automorphic representations organizing Hecke actions and their local parameters \cite{Bump1997}.
\end{itemize}
Motives provide a conjecturally universal organization principle for these compatibilities, with periods appearing as comparison invariants between realizations.

The present work motivates the following objective for the audit program:
\begin{equation}
\label{eq:diagram_objective}
\ScanAlg
\xrightarrow{\ \HSP\ }
\PerDatum
\xrightarrow{\ \mathrm{Mot}\ }
\Cat{Mot}
\xrightarrow{\ \mathrm{Real}_\ell\ }
\Cat{GalRep}
\qquad\text{and}\qquad
\Cat{Mot}
\xrightarrow{\ \mathrm{Aut}\ }
\Cat{AutRep},
\end{equation}
where:
\begin{itemize}
\item $\Cat{Mot}$ is a suitable category of motives (or a computable subcategory, e.g.\ mixed Tate motives),
\item $\mathrm{Mot}$ attaches a motive to a period datum when available (programmatic),
\item $\mathrm{Real}_\ell$ takes an $\ell$-adic realization, producing a Galois representation,
\item $\mathrm{Aut}$ associates an automorphic representation (or parameter) in regimes where Langlands correspondences are known.
\end{itemize}

In this diagram, the arrow $\HSP$ is \emph{closed} and auditable within this paper. The arrows $\mathrm{Mot}$, $\mathrm{Real}_\ell$, and $\mathrm{Aut}$ are classical in several standard regimes (e.g.\ mixed Tate motives and modular motives), but they are not available as a uniform construction on arbitrary period data. In the present paper we therefore treat them as programmatic interfaces and aim to instantiate them only in controlled test cases.

Appendix~\ref{app:motivic_worked_examples} records two such classical test cases. The first is the mixed Tate period calculus on $\mathbb{P}^1\setminus\{0,1,\infty\}$, where iterated-integral periods include $\log 2$ and $\zeta(2)$. The second is the modular/elliptic setting, where weight-$2$ newforms have period integrals controlling $L(f,1)$ and admit $\ell$-adic realizations with Frobenius traces given by Hecke eigenvalues.

\subsection{How this relates to the stairway chain}
\label{sec:relation_to_stairway_chain}

The stairway chain \cite{Ma2025STI} produces, within Layer~0/1, a route from scan protocols to the Hecke prime skeleton and cusp coefficient data. In standard arithmetic geometry, the same coefficient data can be interpreted as part of the realization data of modular motives: Hecke eigenvalues correspond to Frobenius traces in suitable $\ell$-adic realizations, while periods arise from comparison between Betti and de~Rham realizations.

The contribution of the present paper is to isolate a stable numerical interface in the middle: period data that are directly computable from scan averages. This isolates an auditable passage
\[
\text{protocol}\ \Rightarrow\ \text{period datum}\ \Rightarrow\ \text{(candidate motive source)},
\]
and clarifies the division of labor:
\begin{itemize}
\item the \emph{closed} part explains why periods (hence $\pi,\log,\zeta$-values) emerge from protocols and how finite-resource errors are controlled;
\item the \emph{programmatic} part aims to integrate this period interface with the Hecke/Langlands interface of \cite{Ma2025STI} to obtain a genuinely functorial closure.
\end{itemize}

\subsection{A controlled modular-motive test case}
\label{sec:controlled_modular_motive_case}

As a concrete target for future work, one may restrict to modular forms of low weight/level and the associated modular motives. In this regime, both sides admit explicit computations: cusp forms have Fourier expansions with Hecke eigenvalues; their periods can be computed as integrals of modular forms on suitable cycles; and Deligne's construction attaches $\ell$-adic Galois representations with matching traces \cite{Deligne1971Bourbaki,Deligne1974}. Appendix~\ref{app:motivic_worked_examples} spells out this classical chain at the level of period integrals and Frobenius traces.

The scan-period interface of this paper suggests a computable experiment: choose kernels whose associated period data are known to coincide with modular periods, and compare their scan realizations with the cusp/Hecke data extracted by the stairway pipeline. Such a comparison would provide an auditable bridge from scan protocols to motivic realizations in a controlled modular setting.


